
\chapter{Evaluation}
\label{cha:evaluation}

This chapter evaluates the implementation of the reasoning layer described in \cref{cha:implementation}. Consistent with the scope defined in Section~\ref{sec:purpose-goals}, it is a functional evaluation of the reasoning layer, not an external validation of assembly correctness by domain experts. It tests whether the implemented rules, validation logic, confidence handling, dependency tracking, and graph export behave consistently with the design specified in \cref{cha:design}. Therefore, the reported validation statuses should be interpreted as rule-based validation outcomes, not as independently verified judgments of whether each physical assembly action was correct.

A negative result in this evaluation would occur if the reasoning layer failed to preserve the input order, failed to generate required artifacts, inferred constraints without traceable rule provenance, allowed rejected or invalidated effects to support later dependencies, failed to mark unsupported requirements as missing, failed to downgrade low-confidence evidence, or exported graph relations that contradicted the validation records. These failures would show that the proposed reasoning representation does not reliably support traceable validation, even under the assumed symbolic input conditions.

This evaluation does not compare the rule base against expert assembly judgment; such an external semantic validation is outside of the implemented evaluation.


\section{Evaluation objective}
\label{sec:evaluation-objective}

The objective of the evaluation is to determine whether the implemented reasoning
layer can transform ordered step artifacts into traceable procedural reasoning
outputs. The evaluation therefore examines five aspects: artifact correctness,
constraint inference coverage, validation behavior under uncertain or conflicting
evidence, graph-based traceability, and robustness under controlled symbolic
input degradation.

The alignment with the research questions is as follows. RQ1 is addressed by evaluating whether ordered step hypotheses can be represented through step records, predicates, constraints, validation records, and graph nodes. RQ2 is addressed by evaluating how confidence values are propagated from predicates to constraints and validation outcomes. RQ3 is addressed through controlled perturbations of the evidence, such as missing requirements, reduced confidence, and incompatibility conditions. RQ4 is addressed by evaluating whether validation outcomes can be traced back to predicates, inferred constraints, rules, missing requirements, incompatibilities, and dependency relations.

\section{Evaluation setup}
\label{sec:evaluation-setup}

The IndustReal-based upstream pipeline provides the external assembly context for
the reasoning-layer evaluation. From it, Layers 1 and 2 produce ordered step and graph-like artifacts from oracle-style symbolic labels and use them as inputs to
the reasoning layer. The results reported here evaluate whether the proposed reasoning layer can validate, enrich, and expose the structure of these
ordered procedural artifacts.

The reasoning-layer evaluation starts from the adapter outputs:
\begin{itemize}
    \item \path{step_records.jsonl}, containing ordered step records;
    \item \path{predicates.jsonl}, containing symbolic predicates associated with steps, objects, component types, tools, and source information.
\end{itemize}

Layer 3 consumes these files and writes:
\begin{itemize}
    \item \path{inferred_constraints.csv}, containing inferred requirements, expected effects, tool requirements, safety requirements, and incompatibilities.
\end{itemize}

Layer 4 consumes the step records, predicates, and inferred constraints, and writes:
\begin{itemize}
    \item \path{validation_records.jsonl};
    \item \path{step_validations.csv};
    \item \path{explanation_traces.json};
    \item \path{effect_history_diagnostics.csv}.
\end{itemize}

Finally, the graph exporter writes:
\begin{itemize}
    \item \path{procedural_reasoning_graph.json};
    \item \path{procedural_reasoning_graph_nodes.csv};
    \item \path{procedural_reasoning_graph_edges.csv}.
\end{itemize}

Recent graph exports also include graph-level provenance in
\path{procedural_reasoning_graph.json}. This provenance records the graph build
time, graph schema version, builder function, input artifact hashes, and the
versions and SHA-256 hashes of \path{config/domain_config.yaml} and
\path{config/thesis_rules.yaml}. When imported into Neo4j, the same provenance
is exposed through a \texttt{GraphManifest} node for the imported
\texttt{graph\_name}. These metadata fields are used only to check freshness and
reproducibility; they do not change the validation rules or the procedural
semantics of the graph.

The evaluation uses these files to check whether the reasoning pipeline behaves according to the design specified in \cref{cha:design}.

\section{Evaluation 1: Pipeline artifact consistency}
\label{sec:evaluation-artifact-correctness}

This evaluation verifies that each pipeline stage can be inspected independently and
that no reasoning stage depends on hidden state.

\paragraph{Goal:}
To check if the implemented pipeline produces the expected artifacts and whether these artifacts satisfy basic consistency conditions. 

\paragraph{Description:}
The evaluation was done on the IndustReal clip \path{03_assy_0_1} in the \path{od_only} mode.

\paragraph{Results:}
The content of the produced artifacts is summarized in \cref{tab:evaluation-artifact-checks}.
To complement the artifact-level checks for the evaluated sample, the input step sequence in \Cref{tab:evaluation1-input-steps} is compared with the generated procedural reasoning graph shown in \Cref{fig:evaluation1-procedural-graph}. The figure focuses on the validated step nodes and on the two relations most relevant to this evaluation: temporal ordering through \texttt{NEXT} edges and inferred procedural dependencies through \texttt{DEPENDS\_ON} edges. 

\begin{table}[htbp]
\centering
\caption{Pipeline artifact checks.}
\label{tab:evaluation-artifact-checks}
\begin{tabular}{p{0.22\textwidth} p{0.34\textwidth} p{0.36\textwidth}}
\toprule
Check & Expected condition & Result \\
\midrule
Step records produced & Each input step is represented in \path{step_records.jsonl}. & 11 step records cover 11 input steps. \\
Predicate records produced & Each step has associated symbolic predicate evidence where available. & 106 predicates reference valid step records. \\
Layer 3 constraints produced & Rule inference writes \path{inferred_constraints.csv}. & 28 constraints produced: 11 \texttt{produces}, 13 \texttt{requires}, 3 \texttt{requiresSafety}, and 1 \texttt{requiresTool}. \\
Layer 4 validation records produced & Each step has a validation record and status. & 11 validation records include statuses. \\
Explanation traces produced & Validation decisions include trace information. & Each validation decision includes trace information. \\
Graph export produced & Graph JSON and node/edge CSV files are written. & Graph export contains 198 nodes and 513 edges. \\
Input order preserved & \texttt{NEXT} edges follow the validation index. & 10 \texttt{NEXT} edges follow validation order. \\
Rejected-step dependency rule respected & Rejected steps do not support later \texttt{DEPENDS\_ON} edges. & 9 \texttt{DEPENDS\_ON} edges avoid rejected-step support. \\
Rule coverage diagnostics & Steps with predicate evidence but no matching Layer 3 rule are explicitly reported. & All 11 steps have rule coverage; no rule-coverage warning was emitted. \\
\bottomrule
\end{tabular}
\end{table}

\paragraph{Graph provenance note:}
The Evaluation~1 sample graph passed the artifact and consistency checks, but
the local graph artifact used for this sample predates the current graph-level
provenance metadata. The Evaluation~1 report therefore marks graph provenance as
unavailable for this sample and recommends rebuilding the graph when freshness
of the domain and rule configuration must be verified. This does not affect the
artifact consistency checks in \Cref{tab:evaluation-artifact-checks}; it only
limits what can be concluded about the exact config hashes used to build that
particular sample graph.

\begin{table}[htbp]
\centering
\caption{Input step sequence used in Evaluation 1.}
\label{tab:evaluation1-input-steps}
\begin{tabular}{p{0.10\textwidth} p{0.18\textwidth} p{0.62\textwidth}}
\toprule
Index & Source frame & Input step description \\
\midrule
0 & 709 & Install base \\
1 & 709 & Install rear chassis \\
2 & 709 & Install front rear chassis pin \\
3 & 1187 & Install front chassis \\
4 & 1187 & Install front chassis pin \\
5 & 1788 & Install rear rear chassis pin \\
6 & 1788 & Install front bracket \\
7 & 1788 & Install front bracket screw \\
8 & 1788 & Install front wheel assy \\
9 & 2735 & Remove front wheel assy \\
10 & 2735 & Install rear wheel assy \\
\bottomrule
\end{tabular}
\end{table}


\paragraph{Analysis:}
For this evaluated sample, the visual inspection supports the artifact checks by showing that the exported graph preserves the input step order and represents dependency links without allowing rejected steps to support later steps.

\begin{figure}[htbp]
\centering
\includegraphics[width=1.0\textwidth]{Figures/Eval01_03_ass_0_1_Step_NEXT_DEPENDS_noErrors.png}
\caption{Procedural reasoning graph generated for Evaluation 1, showing validated step nodes, temporal \texttt{NEXT} edges, inferred \texttt{DEPENDS\_ON} edges, and removal-specific reasoning. The removal action (Step 9) is preserved as a step in the temporal sequence and is represented by explicit remove-action constraints in the current rule set.}
\label{fig:evaluation1-procedural-graph}
\end{figure}

The graph also shows the Layer 4 validation status of each step. Steps marked as \texttt{uncertain} are not graph-export errors; they indicate that the validation stage preserved cases where the available evidence or inferred support was insufficient for full acceptance. In this evaluation, the status labels are used to
check artifact propagation rather than to analyze the correctness of each decision.

\subsection{Rule coverage diagnostics}
\label{sec:evaluation-rule-coverage-diagnostics}

The artifact checks also include a rule-coverage diagnostic for steps that contain symbolic predicate evidence but do not trigger any Layer 3 rule. This distinction is important because the absence of inferred constraints may either mean that a step has no procedural dependency, or that the current rule base does not cover the observed action type.

In the current artifact-correctness run, this diagnostic does not report an uncovered action. Step~9, \textit{Remove front wheel assy}, is preserved in the input sequence, has symbolic evidence, and triggers removal-specific rules that require a prior installed state and produce a removed-state effect. The absence of a warning is therefore itself an artifact-level check: the rule-coverage report confirms that the observed removal action is covered by the current rule configuration.

This diagnostic strengthens traceability by making unsupported actions visible in the generated artifacts. It separates rule-base limitations from validated dependency decisions and prevents observed but unsupported actions from passing through the reasoning chain silently.

\section{Evaluation 2: Constraint inference coverage}
\label{sec:evaluation-constraint-coverage}

The second evaluation checks whether Layer 3 enriches symbolic predicate records
with procedural constraints, regardless of the perception accuracy or the final validation status of each step. 

\paragraph{Goal:}
To verify that the rule layer derives requirements, expected effects, tool requirements, safety requirements, removal effects, and incompatibility constraints from the available symbolic facts. 

\paragraph{Description:}
The evaluation was run over 39 reasoning-layer output folders from IndustReal, covering outputs produced in the \texttt{od\_only} mode, the \texttt{od\_plus\_psr\_error\_hints} mode, and one additional sample output. In the error-hint mode, the upstream oracle pipeline was allowed to use annotations from \texttt{PSR\_labels\_with\_errors.csv} when such annotations were present. Therefore, the mode label does not imply that every corresponding clip contains an error.

\paragraph{Results:}
For readability, \Cref{tab:evaluation-constraint-coverage} reports a representative subset using
two clips, each evaluated both without and with error-hint input. 

\begin{table}[htbp]
\centering
\caption{Constraint inference coverage for selected evaluation clips.}
\label{tab:evaluation-constraint-coverage}
\small
\begin{tabular}{p{0.18\textwidth}crrrrrr}
\toprule
Clip & \makecell{Error hints\\enabled} & Predicates & Requires & Produces & Tools & Safety & Incompat. \\
\midrule
\texttt{03\_assy\_0\_1} & No  & 123 & 16 & 11 & 1 & 6 & 0 \\
\texttt{03\_assy\_0\_1} & Yes & 123 & 16 & 11 & 1 & 6 & 0 \\
\texttt{03\_assy\_1\_3} & No  & 130 & 17 & 12 & 1 & 6 & 0 \\
\texttt{03\_assy\_1\_3} & Yes & 175 & 20 & 14 & 1 & 8 & 2 \\
\bottomrule
\end{tabular}
\end{table}

The selected clips illustrate how the same rule layer responds to different upstream evidence. For \texttt{03\_assy\_0\_1}, enabling error hints does not change the inferred constraint counts because this clip does not contain explicit error events. In contrast, \texttt{03\_assy\_1\_3} contains additional evidence when the
error-hint mode is enabled. This increases the number of predicates and inferred
constraints, including two incompatibility constraints derived from explicit
error actions. To make the coverage numbers interpretable, \Cref{tab:evaluation-constraint-examples} shows representative Layer 3 outputs from \texttt{03\_assy\_1\_3} with error hints
enabled. 

\begin{table}[htbp]
\centering
\caption{Representative Layer 3 constraint inference examples from \texttt{03\_assy\_1\_3} with error hints enabled.}
\label{tab:evaluation-constraint-examples}
\scriptsize
\begin{tabular}{p{0.09\textwidth} p{0.20\textwidth} p{0.40\textwidth} p{0.23\textwidth}}
\toprule
Event & Rule case & Inferred constraints & Interpretation \\
\midrule
0 &
Install base &
\texttt{produces(installed, base, workspace)} &
The base step creates an initial installed-state effect. \\

2 &
Install chassis pin &
\texttt{requires(installed, rear\_chassis, base)}\newline
\texttt{produces(installed, front\_rear\_chassis\_pin, rear\_chassis)}\newline
\texttt{requires(aligned, front\_rear\_chassis\_pin, rear\_chassis)}\newline
\texttt{requiresSafety(secured, base, workspace)} &
One step can generate a precondition, an effect, an assembly condition, and a safety requirement. \\

3 &
Error action &
\texttt{incompatibleAction(front\_chassis, error)} &
An explicit error action is converted into an incompatibility constraint. \\

7 &
Remove rear chassis &
\texttt{requires(installed, rear\_chassis, base)}\newline
\texttt{produces(removed, rear\_chassis, base)} &
Removal is represented as a state-changing action with a required installed state and a removed effect. \\

11 &
Install front bracket screw &
\texttt{requires(installed, front\_bracket, front\_chassis)}\newline
\texttt{produces(installed, front\_bracket\_screw, front\_bracket)}\newline
\texttt{requiresTool(screwdriver)} &
Domain knowledge about screws adds a tool requirement. \\

14 &
Error action &
\texttt{incompatibleAction(rear\_wheel\_assy, error)} &
A second explicit error action is also represented as an incompatibility. \\
\bottomrule
\end{tabular}
\end{table}

\paragraph{Analysis:}
The examples in \Cref{tab:evaluation-constraint-examples} clarify what the
coverage counts represent. Layer 3 does not merely copy upstream step labels; it uses action predicates, component types, installation targets, and domain rules to infer additional procedural structure. Event~2 generates a precondition, an expected effect, an assembly condition, and a safety requirement from one installation step. Events~3 and~14 show how error-hint evidence is converted into
explicit incompatibility constraints, while Event~7 shows how a removal action is represented through both a required installed state and a produced removed state.

Across all evaluated folders, Layer 3 produced 2001 inferred constraints from 7580 predicate records over 670 step records: 654 \texttt{produces}, 918 \texttt{requires}, 68 \texttt{requiresTool}, 345 \texttt{requiresSafety}, and 16 incompatibility constraints. Rule provenance was present for all 2001 inferred constraints. Remove-action semantics were covered for all 150 observed removal steps, with no remove steps reported as \texttt{no\_applicable\_rule}.

This evaluation contributes to RQ1 and RQ2 by showing that symbolic predicates can be transformed into explicit procedural constraints with associated confidence and provenance.


\section{Evaluation 3: Step validation and effect-lifecycle behavior}
\label{sec:evaluation-validation-behavior}

The third evaluation examines Layer 4 behavior. While Evaluation 2 checks whether constraints are inferred, this evaluation checks whether those constraints are used consistently to assign validation outcomes, dependency support, and produced-effect lifecycle states.

\paragraph{Goal:}
The goal is to inspect whether requirements are supported by active previous effects, whether rejected steps are excluded from dependency support, whether removal actions invalidate earlier installed effects, and whether validation statuses are consistent with the configured confidence thresholds and hard incompatibility rules. The evaluation also checks whether lowering supporting confidence can change an otherwise accepted step into an \texttt{uncertain} step without introducing a hard incompatibility.

\paragraph{Description:}
The evaluation uses the IndustReal clip \texttt{08\_assy\_0\_1} in the
\texttt{od\_plus\_psr\_error\_hints} mode. This clip was selected because the
coverage analysis in Evaluation 2 showed that it contains the types of inferred
constraints needed to exercise Layer 4 behavior, including requirements, produced
effects, tool and safety requirements, removal effects, and incompatibility
constraints.

The evaluation focuses on the validation artifacts produced by Layer 4:
\path{validation_records.jsonl}, \path{step_validations.csv},
\path{explanation_traces.json}, and \path{effect_history_diagnostics.csv}. These artifacts are inspected to determine whether validation decisions are
consistent with the inferred constraints and whether the evidence relevant to the
selected validation scenarios is recorded explicitly. A controlled confidence perturbation is also applied to a copy of the reasoning artifacts. This perturbation is used only to test threshold-sensitive behavior.

The validation scenarios are summarized in
\Cref{tab:evaluation-validation-scenarios}.

The distinction between a missing requirement and a dependency is important in this evaluation. A \texttt{requires(...)} constraint expresses a condition needed by a step. Layer 4 checks whether this condition is supported by an earlier active \texttt{produces(...)} effect. If such support is found, the step receives dependency support and the graph may expose it through a \texttt{DEPENDS\_ON} edge. If no active support is found, for example because the producing step was rejected or because the effect was invalidated by a removal action, the
requirement is recorded as missing.


\begin{table}[htbp]
\centering
\caption{Layer 4 validation behaviors evaluated in Evaluation 3.}
\label{tab:evaluation-validation-scenarios}
\begin{tabular}{p{0.26\textwidth} p{0.38\textwidth} p{0.26\textwidth}}
\toprule
Behavior & Validation question & Expected outcome \\
\midrule
Requirement support &
When a step has a \texttt{requires(...)} constraint, does Layer 4 find a matching active earlier \texttt{produces(...)} effect? &
Supported requirements are recorded in dependency support; unsupported requirements are recorded as missing. \\

Hard incompatibility &
When Layer 3 infers an \texttt{incompatibleAction(...)} constraint, does Layer 4 treat it as a hard violation? &
The affected step is marked \texttt{rejected}. \\

Rejected-step isolation &
Can a produced effect from a rejected step support a later requirement? &
Rejected steps remain in the trace, but their produced effects are not active support for later steps. \\

Removal invalidation &
When a step produces \texttt{removed(component, target)}, is the earlier active \texttt{installed(component, target)} effect invalidated? &
The installed effect remains historical but is not available as active support after removal. \\

Reduced confidence &
If selected supporting confidence values are lowered below the acceptance threshold, does the status change? &
The affected step becomes \texttt{uncertain} when evidence remains plausible but insufficient for acceptance. \\
\bottomrule
\end{tabular}
\end{table}


\paragraph{Results:}
The selected clip produced 34 validation records: 14 steps were classified as
\texttt{accepted}, 14 as \texttt{uncertain}, and 6 as \texttt{rejected}. The
effect-lifecycle diagnostics contained 20 active effects, 8 invalidated effects,
and 4 effects marked as \texttt{inactive\_rejected}. These results indicate that
Layer 4 did not treat all produced effects as permanently valid; rejected and
invalidated effects were preserved in the trace but excluded from active support.

The scenario-level results are summarized in
\Cref{tab:evaluation-validation-response}. The table separates naturally observed
validation behavior in the selected clip from the controlled confidence
perturbation.

\begin{table}[htbp]
\centering
\caption{Layer 4 validation response in the selected clip and controlled confidence perturbation.}
\label{tab:evaluation-validation-response}
\small
\begin{tabular}{p{0.22\textwidth} p{0.36\textwidth} p{0.08\textwidth} p{0.20\textwidth}}
\toprule
Scenario & Evidence inspected & Result & Supporting artifact \\
\midrule
Requirement support &
57 requirements inspected; 27 were supported and 30 were recorded as missing. &
PASS &
\path{requirement_support_results.csv} \\

Hard incompatibility &
2 incompatibility constraints inspected; both affected steps were rejected with trace evidence. &
PASS &
\path{incompatibility_results.csv} \\

Rejected-step isolation &
6 rejected steps inspected; no rejected step was used as later dependency support. &
PASS &
\path{rejected_step_isolation_results.csv} \\

Removal invalidation &
10 removal effects inspected; no invalidation inconsistencies were found. &
PASS &
\path{removal_invalidation_results.csv} \\

Reduced confidence &
1 controlled perturbation was applied; an accepted step became \texttt{uncertain} after confidence reduction. &
PASS &
\path{reduced_confidence_results.csv} \\
\bottomrule
\end{tabular}
\end{table}

The results show that the validation layer reacts to different evidence
conditions in the intended way. Requirements are either linked to active support
or recorded as missing; hard incompatibilities cause rejection; rejected steps do
not support later dependencies; removal effects invalidate matching installed
effects; and confidence reduction can change an accepted step into an
\texttt{uncertain} step without introducing a hard incompatibility.

This evaluation contributes to RQ2 and RQ3 by testing whether uncertainty,
incompatibility evidence, missing active support, and state invalidation affect
validation outcomes in the intended way. 


\section{Evaluation 4: Procedural graph traceability}
\label{sec:evaluation-graph-traceability}

The fourth evaluation examines whether the procedural reasoning graph exposes the
traceability required by the thesis objective.

\paragraph{Goal:}
Rather than measuring graph coverage over the full dataset, the goal is to check
whether the validation behavior inspected in Evaluation 3 can also be followed in
the graph by traversing from a step to its predicates, inferred constraints,
dependency support, rule provenance, and produced-effect lifecycle.

\paragraph{Description:}
The evaluation uses the same \texttt{08\_assy\_0\_1} clip in the
\texttt{od\_plus\_psr\_error\_hints} mode as Evaluation 3. This choice makes the
connection between validation behavior and graph traceability explicit: the same
steps, dependencies, rejected actions, removal effects, and invalidated effects
that were inspected in Layer 4 are now inspected in the exported procedural
reasoning graph.

The evaluated graph was imported and reported with the per-clip graph name
\texttt{procedural\_reasoning\_graph::raw\_cad\_dataset\_\_all\_test\_clips::od\_plus\_psr\_error\_hints::test\_p1::08\_assy\_0\_1}.
Its graph provenance recorded graph schema version \texttt{1.0}, a local build
timestamp of \texttt{2026-06-30T12:30:08+02:00}, domain model version
\texttt{1.3.0}, and rule-set version \texttt{1.3.0}. The recorded SHA-256
prefixes were \texttt{1f92ec4657d8} for \path{config/domain_config.yaml} and
\texttt{bfff375aa487} for \path{config/thesis_rules.yaml}. These fields make the
graph used for Evaluation~4 distinguishable from stale Neo4j imports with older
or sanitized graph names.

The procedural reasoning graph is evaluated as an explanation and inspection
artifact. The graph should preserve the input order through \path{NEXT} edges
and expose reasoning-derived relations through edges such as \path{REQUIRES},
\path{PRODUCES}, \path{DEPENDS_ON}, \path{HAS_PREDICATE},
\path{HAS_CONSTRAINT}, \path{SUPPORTED_BY}, and \path{DERIVED_FROM}.
For removal actions, the graph should also expose effect lifecycle information
through produced-effect constraint nodes and \path{INVALIDATED_BY} relations.

The evaluated graph contained 34 step nodes, 367 predicate nodes, 91 constraint
nodes, 8 rule nodes, 13 entity nodes, and 47 source nodes. The exported relations
included 33 \path{NEXT} edges, 27 \path{DEPENDS_ON} edges, 57
\path{REQUIRES} edges, 32 \path{PRODUCES} edges, 367 \path{HAS_PREDICATE}
edges, 91 \path{HAS_CONSTRAINT} edges, 458 \path{DERIVED_FROM} edges, and
8 \path{INVALIDATED_BY} edges. The step status distribution was the
same as in Evaluation 3: 14 accepted steps, 14 uncertain steps, and 6 rejected
steps.

\begin{table}[htbp]
\centering
\caption{Procedural reasoning graph traceability checks.}
\label{tab:evaluation-graph-traceability}
\small
\begin{tabular}{p{0.28\textwidth} p{0.36\textwidth} p{0.28\textwidth}}
\toprule
Check & Expected condition & Result \\
\midrule
Order preservation &
Each \texttt{NEXT} edge follows the input validation index. &
PASS: 33 \texttt{NEXT} edges checked; 0 failures. \\

Dependency grounding &
Each \texttt{DEPENDS\_ON} edge is supported by a requirement and an earlier produced effect. &
PASS: 27 dependency edges checked; 0 failures. \\

Requirement visibility &
Requirements are represented as explicit constraint nodes or edges. &
PASS: 71 requirements checked; 0 failures. \\

Missing requirement visibility &
Missing requirements are represented in the validation record and trace. &
PASS: 30 missing requirements visible as graph constraints. \\

Evidence traceability &
Step statuses can be traced to predicates and constraints. &
PASS: 34 step nodes checked; 0 failures. \\

Rule provenance &
Derived constraints can be linked to the rule that produced them. &
PASS: 91 constraints checked; 0 failures. \\

Rejected-step isolation &
Rejected steps do not support later dependency edges. &
PASS: 31 rejected-step graph checks; 0 failures. \\

Provisional dependency visibility &
Dependencies supported by uncertain prior steps are marked as provisional. &
PASS: 4 uncertain-support dependencies checked; 0 failures. \\

Effect invalidation visibility &
Invalidated produced effects are linked to the later effect that invalidated them through \texttt{INVALIDATED\_BY}. &
PASS: 8 invalidated effects checked; 0 failures. \\
\bottomrule
\end{tabular}
\end{table}

\paragraph{Results:}
The results show that the procedural reasoning graph makes the validation
decision inspectable rather than only storing a final status label. Each evaluated
traceability check passed. The graph preserved temporal order, grounded
dependency links in requirement and produced-effect evidence, exposed requirement
and missing-requirement information, connected derived constraints to rule
provenance, prevented rejected steps from supporting later dependencies, marked
uncertain dependency support as provisional, and represented effect invalidation
through \texttt{INVALIDATED\_BY} relations.

The Neo4j views in
\Cref{fig:eval4-step-temporal-order,fig:eval4-step-dependencies,fig:eval4-constraint-evidence,fig:eval4-effect-invalidation,fig:eval4-representative-trace}
provide complementary visual inspections of the same graph. Step node colors
indicate Layer 4 validation status: green denotes \texttt{accepted}, yellow
denotes \texttt{uncertain}, and red denotes \texttt{rejected}. The views separate
the graph into temporal order, dependency support, constraint evidence,
effect-lifecycle relations, and local explanation neighborhoods, since the
complete procedural reasoning graph is too dense to inspect clearly in a single
visualization.

This evaluation contributes to RQ4 by checking whether the graph supports explanation-oriented inspection. It also supports RQ1 by showing that the procedural representation preserves temporal order and reasoning-derived relations, rather than reducing the procedure to a flat sequence of labels. 


\section{Evaluation 5: Symbolic robustness under controlled degradation}
\label{sec:evaluation-symbolic-degradation}

The fifth evaluation examines whether the reasoning layer behaves conservatively
when its symbolic inputs are deliberately degraded.

\paragraph{Goal:}
The goal is to check whether incomplete, low-confidence, contradictory, or
semantically corrupted symbolic evidence changes validation outcomes in a
traceable and conservative way. 

\paragraph{Description:}
The evaluation uses the same \texttt{08\_assy\_0\_1} clip in the
\path{od_plus_psr_error_hints} mode as Evaluations 3 and 4. The baseline validation distribution contained 14 accepted steps, 14 uncertain steps, and 6 rejected steps.

Five controlled perturbations were applied to symbolic inputs before rerunning
the validation layer. The perturbations lower predicate confidence, remove
required support evidence, replace an object reference with an incompatible
object, inject an explicit error action, and remove a produced effect used by
later steps. The expected behavior is conservative degradation: an accepted step
should become \texttt{uncertain} or \texttt{rejected} when the evidence required
for acceptance is weakened, removed, or contradicted. In addition, the perturbed
outputs should preserve explanation traces and diagnostics so that the changed
decision can be inspected.

\begin{table}[htbp]
\centering
\caption{Symbolic input degradation checks.}
\label{tab:evaluation-symbolic-degradation}
\small
\begin{tabular}{p{0.38\textwidth} p{0.52\textwidth}}
\toprule
Perturbation & Result \\
\midrule
Lower predicate confidence &
PASS: the target step changed from \texttt{accepted} to \texttt{uncertain}. \\

Remove required support evidence &
PASS: the target step changed from \texttt{accepted} to \texttt{rejected}; missing requirement evidence was recorded and dependency support was removed. \\

Replace object with incompatible type &
PASS: the target step changed from \texttt{accepted} to \texttt{uncertain}. \\

Inject explicit error action &
PASS: the target step changed from \texttt{accepted} to \texttt{rejected}; no later dependency was supported by the rejected step. \\

Remove produced effect used by later steps &
PASS: seven dependent steps were affected; support was removed, missing requirements were recorded, and no rejected-support violations occurred. \\
\bottomrule
\end{tabular}
\end{table}



\paragraph{Results:}
All five perturbation scenarios passed. Each directly affected accepted step
changed to either \texttt{uncertain} or \texttt{rejected}, meaning that no
accepted decision remained accepted after its supporting evidence was directly
degraded. The confidence degradation scenario changed the target step from
\texttt{accepted} to \texttt{uncertain}. Removing required support evidence,
injecting an explicit error action, and removing a produced effect changed the
affected target steps from \texttt{accepted} to \texttt{rejected}. Replacing an
object reference with an incompatible object changed the affected step from
\texttt{accepted} to \texttt{uncertain}.

The strongest propagation effect occurred when a produced effect used by later
steps was removed. In that scenario, seven dependent steps were affected; each
changed from either \texttt{accepted} or \texttt{uncertain} to
\texttt{rejected}. For each affected dependent step, the dependency support was
removed and missing requirements were recorded. Across all perturbations, no
rejected step supported later dependency edges, and explanation traces were
preserved for all executed scenarios.

The results show that the reasoning layer does not merely keep the original
validation labels when symbolic evidence is degraded. Instead, the validation
outcomes degrade conservatively, and the generated diagnostics make the changed
decisions inspectable. This supports the claim that the reasoning layer can
handle controlled symbolic degradation by exposing low-confidence evidence,
missing requirements, incompatibility evidence, removed dependency support, and
effect-lifecycle consequences.

This evaluation contributes mainly to RQ3 by testing how the reasoning layer
responds when symbolic input quality is deliberately degraded. It also supports
RQ4 because the perturbed outcomes remain traceable through explanation records.



\section{Evaluation Summary}
\label{sec:evaluation-summary}

The evaluation focuses on the reasoning contribution of this thesis. It checks that the implemented pipeline produces the expected artifacts, that Layer 3 enriches symbolic predicates with procedural constraints, that Layer 4 responds according to the specified validation rules when evidence is missing, uncertain, or incompatible, and that the final graph exposes validation status, dependencies, evidence, and provenance. The updated graph export also records graph-level provenance and, after Neo4j import, exposes it through \texttt{GraphManifest} so that reports can distinguish current graphs from stale imports. This evaluation supports the academic claim that the thesis contributes a traceable reasoning layer for validating and enriching ordered procedural artifacts, rather than a perception system or a complete XR guidance application.

\paragraph{Evaluation limitations.}
The evaluations focus on reasoning-layer behavior given upstream symbolic artifacts. They do not evaluate dataset-wide generalization of every graph traceability pattern. Broader limitations are discussed in \Cref{sec:limitations}



\begin{figure}[htbp]
\centering
\includegraphics[width=1.0\textwidth]{Figures/eval04_A_step_status_temporal_order.png}
\caption{Procedural Step nodes ordered by \texttt{NEXT} edges. Node colors indicate validation status as described in the text.}
\label{fig:eval4-step-temporal-order}
\end{figure}

\begin{figure}[htbp]
\centering
\includegraphics[width=1.0\textwidth]{Figures/eval04_B_step_dependencies.png}
\caption{Step sequence with \texttt{DEPENDS\_ON} links exposing dependency support between validated steps.}
\label{fig:eval4-step-dependencies}
\end{figure}

\begin{figure}[htbp]
\centering
\includegraphics[width=0.95\textwidth]{Figures/eval04_C_constraint_evidence -reduced.png}
\caption{Graph view of Step-to-Constraint evidence, including a sample of requirements, produced effects, and incompatibilities. Step colors indicate validation status; other node colors follow the Neo4j node-type styling.}
\label{fig:eval4-constraint-evidence}
\end{figure}

\begin{figure}[htbp]
\centering
\includegraphics[width=1.0\textwidth]{Figures/eval04_D_effect_lifecycle_invalidation.png}
\caption{A sample of produced-effect constraints linked by \texttt{INVALIDATED\_BY} relations to later removal effects.}
\label{fig:eval4-effect-invalidation}
\end{figure}

\begin{figure}[htbp]
\centering
\includegraphics[width=1.0\textwidth]{Figures/eval04_E_full_representative_trace_step17.png}
\caption{Compact explanation neighborhood for a representative Step, showing evidence and provenance links. Step colors indicate validation status; other node colors follow the Neo4j node-type styling.}
\label{fig:eval4-representative-trace}
\end{figure}
